%! AP Statistics — Unit 4, Lesson 4.3 — Warm-Up KEY
\documentclass[10pt]{article}
\usepackage{apstats-article}
\usepackage{apstats-key}

\begin{document}

\pageheader{Unit 4, Lesson 4.3}{Warm-Up --- KEY}
\namedateperiod

\vspace{0.3in}

{\large\bfseries\color{navy} Part 1 \quad Spiral Review}

\vspace{0.12in}

\begin{enumerate}

    \item Steps to design a simulation:
    \begin{enumerate}[label=(\arabic*), itemsep=6pt]
        \item \ans{Identify the random process to be modeled.}
        \item \ans{Assign values from your chance device to each possible outcome.}
        \item \ans{Choose an appropriate chance device whose probabilities match the real-world process.}
        \item \ans{Run trials and record the outcome of each trial systematically.}
        \item \ans{Compute the relative frequency: (number of times event occurred) $\div$ (total number of trials).}
    \end{enumerate}

    \vspace{0.16in}

    \item
    \begin{enumerate}[label=(\alph*), itemsep=6pt]
        \item \ans{$\hat{p} = 104/200 = 0.52$}
        \item \ans{Yes. The Law of Large Numbers states that as the number of trials increases, the empirical probability tends to get closer to the true probability. With 2{,}000 trials, random fluctuations average out more and the estimate would likely be closer to 0.49 than the current 0.52.}
    \end{enumerate}

\end{enumerate}

\vspace{0.25in}

{\large\bfseries\color{navy} Part 2 \quad Looking Ahead}

\vspace{0.12in}

\noindent Today we derive exact (theoretical) probabilities --- no simulation needed.

\vspace{0.14in}

\begin{tcolorbox}[colback=greenbg, colframe=greenacc, boxrule=0.8pt, arc=2mm,
    left=4mm, right=4mm, top=3mm, bottom=3mm]
    \small
    \begin{itemize}[leftmargin=1.3em, itemsep=8pt]
  \item The \textbf{sample space} lists every possible outcome; no outcome is counted twice.
  \item When all outcomes are equally likely: $P(E) = \dfrac{\text{outcomes in }E}{\text{total outcomes}}$.
  \item The complement rule: $P(\text{not }E) = 1 - P(E)$.
    \end{itemize}
\end{tcolorbox}

\vspace{0.2in}

\noindent\textcolor{slate}{\small Estimate: probability of rolling 3 or higher on a fair die.}\\[8pt]
\ans{The sample space is \{1, 2, 3, 4, 5, 6\}. Outcomes $\geq 3$ are \{3, 4, 5, 6\}, so $P(\geq 3) = 4/6 \approx 0.667$.}

\end{document}
